\documentclass[10pt,a4paper]{amsart}
\usepackage[margin=19mm]{geometry}
\usepackage{amsmath,amssymb,mathtools}
\usepackage[scr=rsfs,cal=euler]{mathalfa}
\usepackage{xfrac}
\usepackage[unicode,hidelinks,bookmarks=false]{hyperref}
\newcommand{\ie}{i.e.\ }
\newcommand{\cf}{cf.\ }
\newcommand{\resp}{respectively\ }
\newcommand{\Subsection}{\subsection}
\providecommand{\nobreakdash}{\nobreak\hbox{-}}
\setlength{\emergencystretch}{2em}
\multlinegap0pt
\usepackage{bm}
\usepackage{bbm}
\usepackage{dsfont}
\usepackage{xspace}
\usepackage{cancel}
\usepackage{mathtools}
\usepackage{amsthm}
\makeatletter
\@ifundefined{thm}{\newtheorem{thm}{Theorem}}{}
\@ifundefined{core}{\newtheorem{core}[thm]{Corollary}}{}
\@ifundefined{lem}{\newtheorem{lem}[thm]{Lemma}}{}
\makeatother
\def \geq {\geqslant}
\def \leq {\leqslant}
\title[On $2$-integral Cayley graphs]{On $2$-integral Cayley graphs}
\author{Alireza Abdollahi, Majid Arezoomand, Tao Feng, Shixin Wang}
\date{Source: arXiv:2401.15306v3 (CC BY 4.0)}
\begin{document}
\maketitle

\noindent\emph{Source and licence.} On $2$-integral Cayley graphs. Version arXiv:2401.15306v3 (CC BY 4.0), distributed under the Creative Commons Attribution 4.0 International licence, as recorded in the arXiv licence metadata for this exact version and shown on the abstract page https://arxiv.org/abs/2401.15306v3. Full rights notice: licenses/arxiv-2401.15306-CC-BY-4.0.txt.\par\smallskip

\noindent\textit{Editorial scope.} Counted unit: the Thm whose source label is \texttt{4} in the article (source0.tex lines 657--675, source label \texttt{4}), delivered together with the article's own complete original proof of it (source0.tex lines 676--724, one \texttt{proof} environment of 49 lines). No mathematics is written, repaired or re-derived: statement and proof are the article's own text line for line. Required context retained uncounted: half (lines 527--530); asli (lines 555--563); 2 (lines 588--591); +1 (lines 601--603); 3 (lines 610--617); induced (lines 640--648). Every definition, display and label of those lines is retained unchanged. Omitted from this package and not redistributed: 1--526, 531--554, 564--587, 592--600, 604--609, 618--639, 649--656, 725--1148. Nothing inside a retained excerpt is omitted or altered: every retained body is delivered exactly as the article's lines are written, so no omission receipt of any kind accompanies this package. Citations: the retained excerpts carry no citation command at all, so no bibliography entry of the article is transcribed and no reference is redistributed. Reference adaptation: none was needed and none was made; every reference inside the delivered unit resolves inside it, so no unresolved reference or citation is produced. Printed slips kept literal: no printed word, symbol, hypothesis or number of the counted statement or of its proof was altered. Layout and class: the article's own class is \texttt{article} with packages including latexsym, amssymb, amsfonts, amsmath, amsthm, nccmath, float, enumitem, setspace, bm, authblk, xcolor, hyperref, geometry; the wrapper is the lane's pinned \texttt{amsart} editorial wrapper at 10pt on A4, which restates the theorem-like environments the retained text needs and adds the article's own macro definitions that the retained text needs (\texttt{\textbackslash geq}, \texttt{\textbackslash leq}), with extra wrapper packages (bm, bbm, dsfont, xspace, cancel, mathtools). The delivered numbering is the unit's own. Assets: no retained span contains a figure environment, a graphics-inclusion command, a PDF-page inclusion, a table or a TikZ picture; no image file is copied into this package, the package declares no asset dependency and no image file is redistributed with it. Licence: version arXiv:2401.15306v3 of this article is distributed under the Creative Commons Attribution 4.0 International licence, as recorded in the arXiv licence metadata for this exact version and shown on the abstract page https://arxiv.org/abs/2401.15306v3. A full rights notice is delivered at licenses/arxiv-2401.15306-CC-BY-4.0.txt. Characters: every retained excerpt is pure ASCII, so the delivered engine, XeLaTeX with its default Latin Modern text font, reports no missing character.
	\begin{lem}\label{half}
		Let $G$ be a finite  group, $g\in G$ be an element of order $n$, $\varnothing\neq S\subseteq [g]$, maybe not inverse-closed,  and $\Gamma=\mathrm{Cay}(G,S)$. 
		Then $\Gamma$ is $2$-integral if and only if there exists a subgroup $H$ of $\mathrm{Aut}(\langle g\rangle)$ such that $|\mathrm{Aut}(\langle g\rangle):H|=2$, $S=s^H$ for all $s\in S$ and $|S|=\frac{\varphi(n)}{2}$.
	\end{lem}

	\begin{core}\label{asli}
		Let $G$ be a finite group, $\Omega(G)=\{[g_1],\ldots,[g_n]\}$, $S$ is a normal subset of $G$ and $\Gamma=\mathrm{Cay}(G,S)$.
		Then $S=\bigcup_{i=1}^{n}S_i$, where $S_i=[g_i]\cap S$ for each $1\leq i\leq n$.
		If $\Gamma$ is $2$-integral, then for $S_i\neq \varnothing$,
		\begin{itemize}
			\item [$(1)$]  there exists $s_i\in S$ such that $S_i=[s_i]$ or $S_i=s_i^H\subset [s_i]$ for some subgroup $H$ of $\mathrm{Aut}(\langle s_i\rangle)$ of index $2$, and $|S_i|=\frac{\varphi(n_i)}{2}$, where $n_i=o(s_i)$;
			\item[$(2)$] there exists $1\leq i_0\leq n$ such that $S_{i_0}$ is of the latter form. 
		\end{itemize}
	\end{core}

	\begin{thm}\label{2}
		Let $G$ be a finite abelian group. Then $G\in\mathcal{G}_2$ if and only if $G\cong\mathbb Z_n$, where $n=5,8,10,12$.
		Furthermore, a cycle with order $n$ is $2$-integral if and only if $n=5,8,10,12$.
	\end{thm}

	\begin{lem}\label{+1}
		Let $G$ be a finite group and $S$ be a normal subset of $G$.  Then for any $1\neq x\in Z(G)$ in which $S\cap [x]=\varnothing$, $\mathrm{Cay}(G,S)$ and $\mathrm{Cay}(G,S\cup [x])$ have the same splitting field. In particular, if $x\in Z(G)\setminus S$ is an involution, then $\mathrm{Cay}(G,S)$ and $\mathrm{Cay}(G,S\cup \{x\})$ have the same splitting field.
	\end{lem}

	\begin{thm}\label{3}
		Let $G$ be a finite abelian group. Then $G\in\mathcal{G}_3$ if and only if $G$ is isomorphic to $\mathbb Z_n$ or $\mathbb Z_n\times\mathbb Z_2$, where $n=8,10,12$.
		Furthermore, a connected cubic abelian undirected Cayley graph $\Gamma$ is $2$-integral if and only if it is isomorphic to one of the  following $6$ Cayley graphs $\mathrm{Cay}(G,S)$, where:
		\begin{itemize}
			\item[$\romannumeral1$.] $G=\langle x\rangle \cong \mathbb{Z}_n$ where $n=8,10,12$ and $S=\{x,x^{-1},x^{\frac{n}{2}}\}$;
			\item[$\romannumeral2$.] $G=\langle x\rangle \times \langle y\rangle\cong \mathbb{Z}_n\times \mathbb{Z}_2$ where $n=8,10,12$ and $S=\{x,x^{-1},y\}$.
		\end{itemize}
	\end{thm}

	\begin{lem}\label{induced}
		Let $G$ be a finite group, $S_1,\ldots,S_k$ be normal subsets of $G$, $\langle S_i\rangle\cap\langle S_j\rangle=\{1\}$ for all distinct $i,j$, and $S=S_1\cup\cdots\cup S_k$. Let $\Gamma=\mathrm{Cay}(G,S)$ and  $\Gamma_i=\mathrm{Cay}(G,S_i)$, $i=1,\ldots,k$. Then
		\begin{itemize}
			\item[$(1)$] $\mathbb{SF}(\Gamma_i)\subseteq\mathbb{SF}(\Gamma)$ for all $i$,
			\item[$(2)$] $\deg(\Gamma_i)$ divides $\deg(\Gamma)$ for all $i$,
			\item[$(3)$] if $\deg(\Gamma_i)=\deg(\Gamma)=2$ for some $i$, then $\mathbb{SF}(\Gamma)=\mathbb{SF}(\Gamma_i)$,
			\item[$(4)$] if $\mathbb{SF}(\Gamma_i)=\mathbb{F}$ for all $i$, then $\mathbb{SF}(\Gamma)=\mathbb{F}$.
		\end{itemize}
	\end{lem}

	\begin{thm}\label{4}
		Let $G$ be a finite abelian group. Then $G\in\mathcal{G}_4$ if and only if $G$ is isomorphic to one the following groups:
		\begin{itemize}
			\item[$(1)$] $\mathbb Z_n$, where $n=8,10,12,15,16,20,24,30$.
			\item[$(2)$] $\mathbb Z_n\times\mathbb Z_2^2$, where $n=5,8,10,12$.
			\item[$(3)$] $\mathbb Z_n\times\mathbb Z_2$, where $n=8,10,12$.
			\item[$(4)$] $\mathbb Z_n\times\mathbb Z_m$, where $(n,m)$ is one of the pairs $(3,12)$, $(4,8)$, $(4,10)$, $(4,12), (6,8)$, $(6,10)$, $(6,12)$, $(5,5)$, $(5,10)$, $(8,8)$, $(10,10)$, $(12,12)$.
		\end{itemize}
		Furthermore, a connected $4$-regular abelian undirected Cayley graph $\Gamma$ is $2$-integral if and only if it is isomorphic to one of the  following $39$ Cayley graphs $\mathrm{Cay}(G,S)$, where
		\begin{itemize}
			\item[$\romannumeral1$.] $G=\langle x\rangle \cong \mathbb{Z}_n$ where $n=15,16,20,24,30$, and $S=\{x,x^{-1},x^k,x^{-k}\}$ where $1\leq k\leq n$ such that $(k,n)=1$,
			\item[$\romannumeral2$.] $G=\langle x\rangle \cong \mathbb{Z}_8$ and $S=\{x,x^2,x^6,x^7\}$,		
			\item[$\romannumeral3$.] $G=\langle x\rangle \cong \mathbb{Z}_{10}$ and $S=\{x,x^2,x^8,x^9\}$,
			\item[$\romannumeral4$.] $G=\langle x\rangle \cong \mathbb{Z}_{12}$ and $S=\{x,x^2,x^{10},x^{11}\}$ or $\{x,x^3,x^9,x^{11}\}$ or $\{x,x^4,x^8,x^{11}\}$,
			\item[$\romannumeral5$.] $G=\langle x\rangle\times \langle w\rangle\cong \mathbb{Z}_n\times \mathbb{Z}_2$ where $n=8,10,12$ and $S=\{x,x^{-1},x^{\frac{n}{2}},w\}$,
			\item[$\romannumeral6$.] $G=\langle x\rangle\times \langle z\rangle\times \langle w\rangle\cong \mathbb{Z}_n\times \mathbb{Z}_2^2$ where $n=5,8,10,12$ and $S=\{x,x^{-1},z,w\}$,
			\item[$\romannumeral7$.] $G=\langle x\rangle\times \langle w\rangle\cong \mathbb{Z}_n\times \mathbb{Z}_m$ where $(n,m)$ is one of the pairs $(3,12)$, $(4,8)$, $(4,10)$, $(4,12), (6,8)$, $(6,10)$, $(6,12)$, $(5,5)$, $(5,10)$, $(8,8)$, $(10,10)$, $(12,12)$, and $S=\{x,x^{-1},y,y^{-1}\}$.
		\end{itemize}
	\end{thm}

	\begin{proof}
		Let $S=\{x,y,z,w\}$ be an inverse closed generating set of $G$ and $\mathrm{Cay}(G,S)$  is $2$-integral. Then we deal with the following cases:
		
		\textbf{Case I.} $x^2=y^2=z^2=w^2=1$. In this case, $G$ is an elementary abelian $2$-group isomorphic to $\mathbb Z_2^n$, where $n=3$ or $4$. Furthermore, $S=S_1\cup S_2\cup S_3\cup S_4$,
		where $S_1=\{x\}=[x]$, $S_2=\{y\}=[y]$, $S_3=\{z\}=[z]$ and $S_4=\{w\}=[w]$, which contradicts Corollary \ref{asli}.
		
		\textbf{Case II.} $z^2=w^2=1$, $z,w\notin\langle x\rangle$  and $y=x^{-1}$. In this case, $G=\langle x\rangle\times\langle z\rangle\times\langle w\rangle\cong\mathbb Z_n\times\mathbb{Z}_2^2$, where $n=o(x)$. Furthermore, $S=S_1\cup S_2\cup S_3$, where $S_1=\{x,x^{-1}\}\subseteq [x]$, $S_2=\{z\}=[z]$ and $S_3=\{w\}=[w]$. Now Corollary \ref{asli} implies
		that $S_1\neq [x]$ and $\varphi(n)=4$ which means $n=5,8,10,12$.
		
		\textbf{Case III.} $y=x^{-1}$, $z=x^{\frac{n}{2}}$, where $n=o(x)$ is even and $w^2=1$. In this case, $G\cong\mathbb Z_n\times\mathbb Z_2$. Also $S=S_1\cup S_2\cup S_3$, where
		$S_1=\{x,x^{-1}\}\subseteq [x]$, $S_2=\{x^\frac{n}{2}\}=[x^\frac{n}{2}]$ and $S_3=\{w\}=[w]$. Then by Corollary \ref{asli}, $S_1\neq [x]$ and $\varphi(n)=4$, which means $n=8,10,12$.
		
		\textbf{Case IV.} $y=x^k$ for some $k\neq 1,-1$, $z=x^{-1}$ and $w=y^{-1}$. In this case, $y^2\neq 1$ and $G=\langle x\rangle\cong\mathbb Z_n$, where $n=o(x)$. 
		
		First assume that $(k,n)=1$. Then $S\subset [x]$ and Lemma \ref{half} implies that $\varphi(n)=8$, which means $n=15,16,20,24,30$. 
		
		Now let $(k,n)=d\neq 1$. Then $S=S_1\cup S_2$, where $S_1=\{x,x^{-1}\}\subseteq [x]$ and $S_2=\{x^k,x^{-k}\}\subseteq [x^k]$. 
		If $S_1=[x]$ and $S_2\neq[x^k]$ then $n=3,4,6$ and $\frac{n}{d}=5,8,10,12$, respectively, which is impossible.
		If $S_1\neq [x]$ and $S_2=[x^k]$, then $n=5,8,10,12$ and $\frac{n}{d}=3,4,6$, respectively, which implies $(n,k)=(8,2)$, $(12,2)$, $(12,3)$ or $(12,4)$.
		If $S_1\neq [x]$ and $S_2\neq [x^k]$ then $n,\frac{n}{d}=5,8,10,12$, which implies $(n,k)=(10,2)$.
		
		
		\textbf{Case V.} $\langle x\rangle\cap\langle y\rangle=1$. Then $G=\langle x\rangle\times\langle y\rangle\cong\mathbb Z_n\times\mathbb Z_m$, where $n=o(x)$ and $m=o(y)$. We may assume that $n\leq m$. In this case, $S=S_1\cup S_2$, where $S_1=\{x,x^{-1}\}\subseteq [x]$ and $S_2=\{y,y^{-1}\}\subseteq [y]$. By Corollary \ref{asli}, $S_1=[x], S_2\neq [y]$ or
		$S_1\neq [x],S_2=[y]$ or $S_1\neq [x],S_2\neq [y]$. By a similar discussion to the above cases, in the first case $n=3,4,6,~m=5,8,10,12$, in the second case $n=5,8,10,12,~m=3,4,6$ and in the later case $n,m=5,8,10,12$.
		The second case is impossible because $n\leq m$. 
		Next we will show that in the last case $(n,m)$ must be $(5,5)$, $(5,10)$, $(8,8)$, $(10,10)$ or $(12,12)$.	
		Let $g\in G$ and $o(g)=k\geq 2$ and $C_k=\mathrm{Cay}(\langle g\rangle,\{g,g^{-1}\})$.
		By an easy computation, we have $\mathbb{SF}(C_{5})=\mathbb{SF}(C_{10})=\mathbb Q[\sqrt{5}]$, $\mathbb{SF}(C_8)=\mathbb Q[\sqrt{2}]$ and $\mathbb{SF}(C_{12})=\mathbb Q(\sqrt{3})$. Hence, by Corollary \ref{asli} and Lemma \ref{induced}, we have the result as desired. Moreover, Cases $\mathrm{ \uppercase\expandafter{\romannumeral1}}$-$\mathrm{\uppercase\expandafter{\romannumeral5} }$ prove one direction.
		
		The proof of converse direction of $G\cong\mathbb Z_n\times\mathbb{Z}_2^2$ where $n=5,8,10,12$, and $G\cong\mathbb Z_n\times\mathbb{Z}_2$ where $n=8,10,12$ are similar to the proof of converse direction of Theorem \ref{3}.
		
		Consider the converse direction of case $G=\langle x\rangle \cong \mathbb Z_n$, where $o(x)=n$ is shown in $(1)$.
		Let $\Sigma_{k}=\mathrm{Cay}(\langle x\rangle,\{x,x^k,x^{-1},x^{-k}\})$. 
		First suppose that $n=15,16,20,24,30$ and $(n,k)=1$. 
		Let $\sigma_k:x\mapsto x^k$ and $\tau:x\mapsto x^{-1}$.
		Then $H=\langle \tau,\sigma_k\rangle$ is a subgroup of index $2$ in $\mathrm{Aut}(G)$ and $x^H=\{x,x^k,x^{-1},x^{-k}\}$.
		By Lemmas \ref{half}, we have $\mathrm{SF}(\Sigma_{k})=2$.
		Now suppose that $n=8,10,12$.
		By a tedious computation, one can see that $\mathbb{SF}(\Sigma_{k})=\mathbb Q(\sqrt{2})$ if $(n,k)=(8,2)$, $(12,2)$ or $(12,4)$,
		$\mathbb{SF}(\Sigma_{k})=\mathbb Q(\sqrt{3})$ or $\mathbb Q(\sqrt{5})$ if $(n,k)=(12,3)$ or $(10,2)$ respectively.	
		
		Finally, consider the converse direction of case $G=\langle x\rangle \times \langle y\rangle \cong \mathbb Z_n\times\mathbb Z_m$, where $(n,m)$ shown in $(3)$.
		Let $S=S_1\cup S_2$ and $S_1=\{x,x^{-1}\}\subseteq [x]$ and $S_2=\{y,y^{-1}\}\subseteq [y]$.
		Note that $C_3,C_4,C_6$ are integral with valency $2$ and $\mathbb{Z}_{8},\mathbb{Z}_{10},\mathbb{Z}_{12}$ are all in $\mathcal{G}_2$ by Theorem \ref{2}.
		Hence by Lemma \ref{+1}, $\mathbb Z_n\times\mathbb Z_m \in \mathcal{G}_4$, where $(n,m)=(3,12)$, $(4,8)$, $(4,10)$, $(4,12)$, $(6,8)$, $(6,10)$ and $(6,12)$.
		On the other hand, since $\mathbb{SF}(C_{5})=\mathbb{SF}(C_{10})=\mathbb Q[\sqrt{5}]$, $\mathbb{SF}(C_8)=\mathbb Q[\sqrt{2}]$ and $\mathbb{SF}(C_{12})=\mathbb Q(\sqrt{3})$, 
		by Lemma \ref{induced}, $\mathbb Z_n\times\mathbb Z_m \in \mathcal{G}_4$, where $(n,m)=(5,5)$, $(5,10)$, $(8,8)$. 
		This completes the proof.
	\end{proof}

\end{document}
